%قبل از شروع مسئله ترجیحاً چیزی ننویسید. اما اگر خواستید بنویسید، در قالب لیستی (فقرات و شمارش) بنویسید.
%برای اضافه کردن نمره، بعدی از شروع‌مسئله، [x] را اضافه کنید و به جای x  نمره مورد نظر را قرار دهید. برای زیرمسئله‌ها هم به همین صورت است.
%\شروع‌مسئله[10]
\شروع‌مسئله

فرض کنید که مجموعه کلیدهای ما به صورت بردارهای $[x_{1}, x_{2}, ..., x_{k}]$ باشند، با این فرض که هر کدام از $x_{i}$ها میتوانند از بازه ی $\{0, 1, ..., N - 1\}$ انتخاب شده باشند که $N$ یک مقدار اول است. نشان دهید که خانواده زیر برای این تابع یک خانواده جهانی است:

\begin{flushleft}
	$H = \{h|h(x) = (r_{1}x_{1} + r_{2}x_{2} + ... + r_{k}x_{k}) \% N\}$
\end{flushleft}

که در آن تمام $r_{i}$ ها از بازه $\{0, 1, ..., N - 1\}$ انتخاب شده‌اند.
\solution{
		اگر $H$ یک خانواده از توابع درهم سازی باشد، در این صورت $H$ یک خانواده جهانی است اگر:
	
	\begin{flushleft}
			$\forall x, y \in U, x \neq  y  \qquad \Pr_{h \in H} [h(x) = h(y)] \leq \frac{1}{N}$
	\end{flushleft}
که در آن $N$ اندازه آرایه است. برای محاسبه ی احتمال هم به روش زیر عمل میکنیم:

\begin{flushleft}
	Pr = $\frac{favourable events}{total events}$
\end{flushleft}
 که صورت همان حالت های تساوی $h(x)$ و $h(y)$ است. تعداد حالت های کل، تمام توابع ممکن در مجموعه هستند.  
 
 با توجه به تعریف مجموعه، $r_{i}$ میتواند از بازه ی 0 تا $N - 1$ انتخاب شود. چون تعداد r ها k تا است و هر کدام N مقدار مختلف میتوانند بگیرند، پس $N^{k}$ تا تابع داریم که همان مخرج کسر هستند.
 
 برای محاسبه‌ی تعداد حالت های پیشامد دو کلید مختلف $x = [x_{1}, x_{2}, ..., x_{k}]$ و $y = [y_{1}, y_{2}, ..., y_{k}]$ را در نظر میگیریم. چون متفاوت هستند پس حداقل یکی از مولفه هایشان با هم متفاوت هستند. مثلا فرض کنید مولفه jام این دو با هم متفاوت هستند.
 
 حالا میخوهیم ببینیم چه تعداد حالت هایی وجود دارد که h(x) برابر با h(y) باشد.
 
 \begin{flushleft}
 $h(x) = h(y)$  \\
 $\sum_{i=1}^{k} r_{i}x_{i} \stackrel{N}{\equiv} \sum_{i=1}^{k} r_{i}y_{i}$ \\
  $\sum_{i=1}^{k} r_{i}x_{i} - \sum_{i=1}^{k} r_{i}y_{i} \stackrel{N}{\equiv} 0 $ \\
    $\sum_{i=1}^{k} r_{i}(x_{i} - y_{i}) \stackrel{N}{\equiv} 0 $ \\
    $r_{j}(x_{j}-y_{j}) + \sum_{i=1, i \neq j}^{k} r_{i}(x_{i} - y_{i})  \stackrel{N}{\equiv} 0 $ \\
   $r_{j}(x_{j}-y_{j})  \stackrel{N}{\equiv} - \sum_{i=1, i \neq j}^{k} r_{i}(x_{i} - y_{i})$ \\
$r_{j} \stackrel{N}{\equiv} - \sum_{i=1, i \neq j}^{k} r_{i}(x_{i} - y_{i})(x_{j}-y_{j})^{-1}$ \\
 \end{flushleft}
حالا باید تعداد حالتهای مختلف را محاسبه کنیم. $r_{i}$ دارای N حالت مختلف است.  $r_{j}$ هم به صورت یکتا توسط عبارت سمت راست تعیین میشود. این عبارت $k-1$ جمله دارد. پس به طور کلی $N^{k-1}$ حالت مختلف وجود دارد و:

\begin{flushleft}
	$Pr = \frac{N^{k-1}}{N^{k}}$
\end{flushleft}
پس حکم برقرار است.
}